under movements and deformations of scatterers follows from
\cite{DemersZhang2013,StenlundYoungZhang2013}.  Since $J_R$ is bounded and
piecewise constant on smooth continuity components, multiplication by
$e^{qJ_R}$ is an analytic bounded perturbation.  Kato theory yields the simple
analytic eigenvalue.  The arithmetic local-limit extension follows from the
Young-tower spectral method of \cite{SzaszVarju2004,Gouezel2005}.
\end{proof}

Let $h_{R,q}$ and $\ell_{R,q}$ be normalized right and left leading
eigenvectors.  Their pairing defines an invariant probability
$\mu_{R,q}$ on the natural extension; at $q=0$ it is the Liouville law.

\section{Current pressure, LDP, and endogenous theta}\label{sec:ldp}
\begin{theorem}[Local current LDP and strict convexity]\label{thm:ldp}
For every $R\in\mathcal I$,
\[
 \Lambda_R(q)
 =\lim_{n\to\infty}\frac1n
 \log\int e^{qS_nJ_R}\,d\mu
\]
for $\abs q<q_*$.  It is analytic,
\[
 \Lambda_R'(q)=\int J_R\,d\mu_{R,q},
\]
and
\[
 \Lambda_R''(q)
 =\lim_{n\to\infty}\frac1n
 \Var_{\mu_{R,q}}(S_nJ_R)>0.
\]
After reducing $q_*$ if necessary, $\Lambda_R'$ is a diffeomorphism onto an
open interval $\mathcal J_R$ containing zero.  For $j\in\mathcal J_R$ define
\[
 \theta_R(j)=(\Lambda_R')^{-1}(j).
\]
The local rate function
\[
 I_R(j)=\sup_{\abs q<q_*}\{qj-\Lambda_R(q)\}
\]
satisfies $I_R'(j)=\theta_R(j)$.
\end{theorem}
\begin{proof}
The spectral decomposition gives the cumulant limit and analytic derivatives.
At $q=0$, the planar Lorentz local-limit theorem gives a nondegenerate
covariance matrix for the full cell-displacement cocycle; in particular the
first-coordinate variance is strictly positive.  The eigenvalue Hessian is
continuous in $(R,q)$.  Compactness of the radius interval therefore permits
shrinking the common spectral window so that
$\Lambda_R''(q)>0$ throughout it.  Strict convexity and Legendre duality give
the rest.
\end{proof}

\section{A central-window Gibbs-conditioning theorem}\label{sec:conditioning}
For $n\ge1$ write
\[
 S_{-n,n}J_R=\sum_{k=-n}^{n-1}J_R\circ T_R^k.
\]

\begin{theorem}[Two-sided Liouville Gibbs conditioning]\label{thm:gibbs-conditioning}
Fix compact $\mathcal K\Subset\{(R,q):R\in\mathcal I,\abs q<q_*\}$.
Let $j=\Lambda_R'(q)$ and let integers $k_n$ satisfy
$k_n-2nj=O(1)$.  If $A$ is measurable with respect to collision coordinates
from times $-\ell$ through $\ell$, then
\[
 \left|
 \mu\bigl(A\mid S_{-n,n}J_R=k_n\bigr)-\mu_{R,q}(A)
 \right|
 \le \frac{C_{\mathcal K,\ell}}{\sqrt{n-\ell}}
\]
whenever the conditioning event is admissible.
\end{theorem}
\begin{proof}
Tilt the numerator and denominator by
$e^{qS_{-n,n}J_R-2n\Lambda_R(q)}$.  Fourier inversion for the lattice variable
reduces both to integrals of the complex twisted operators
$\cL_{R,q+it}$.  On $\abs t\le n^{-2/5}$, the simple eigenvalue has the
uniform expansion
\[
 \Lambda_R(q+it)=\Lambda_R(q)+itj-\tfrac12t^2\Lambda_R''(q)+O(t^3),
\]
while the left and right spectral factors propagate for at least $n-\ell$
steps on both sides of the central observable.  Their normalized product is
exactly $\mu_{R,q}(A)$.  Outside the central Fourier window, lattice
aperiodicity and the spectral gap give an exponentially small contribution.
The span-one arithmetic statement is part of the Lorentz local-limit packet and is also visible from the open one-step branches $J_R=0$ and $J_R=1$ in \cref{prop:span}.
Dividing the two local-limit expansions yields the stated rate.  The use of a
central window is essential: a one-sided endpoint observable would retain an
unbalanced eigenfunction factor.
\end{proof}

\section{The driven preparation is the same deterministic billiard}\label{sec:doob}
On the one-sided Young-tower quotient, define
\[
 \cP_{R,q}f
 =e^{-\Lambda_R(q)}h_{R,q}^{-1}\cL_{R,q}(h_{R,q}f).
\]
Then $\cP_{R,q}1=1$ and its stationary law projects to $\mu_{R,q}$.

\begin{proposition}[No microscopic noise]\label{prop:no-noise}
The Markov kernel associated with $\cP_{R,q}$ is the conditional
disintegration of the two-sided invariant measure $\mu_{R,q}$.  Sampling a
full collision state from $\mu_{R,q}$ and applying $T_R$ produces a completely
deterministic trajectory of the original specular table.  The kernel records
only uncertainty created by forgetting stable/future coordinates.
\end{proposition}
